% SLIDE 3 — Problématique clinique
\begin{frame}{Problématique clinique et choix méthodologique}
  \timing{2:00}
  \begin{columns}[T]
    \column{0.58\textwidth}
    \textbf{Enjeu :} distinguer gliome, méningiome et tumeur hypophysaire sur IRM
    conditionne le traitement~\cite{ref_bioengineering11030266,ref_diagnostics15212692}.
    \vspace{0.4em}

    \textbf{Difficultés ML :}
    \begin{itemize}
      \item Haute dimensionnalité ($16{,}384$ pixels bruts)
      \item Différences subtiles de texture et d'intensité
      \item Variabilité d'acquisition ; besoin de métriques robustes (F1-macro)
    \end{itemize}
    \vspace{0.3em}
    \begin{block}{Positionnement}
      \small ML classique + features interprétables (CPU, reproductible)
      vs.\ DL (GPU, données massives, boîte noire)~\cite{ref_diagnostics15212692}
    \end{block}
    \column{0.38\textwidth}
    \centering
    \includegraphics[width=\linewidth]{dataset}
  \end{columns}
\end{frame}

\note{
Le type tumoral dicte le traitement. La lecture manuelle est coûteuse et variable.
Nous adoptons volontairement la voie classique : descripteurs explicites + classifieurs connus,
pour un pipeline auditable et peu coûteux en calcul.
}
